% ------------------------------------------------------------------------------
% 6 Results III -- Simulation studies
% ------------------------------------------------------------------------------
\section{Results III: Simulation Studies}\label{sec:res-sim}

\subsection{Study A: covariate adjustment buys power}

Table~\ref{tab:simA} and Figure~\ref{fig:simA} report 1{,}000 simulated
trials per scenario at $n = 322$ and target hazard ratio $0.72$. With
no prognostic covariates (scaling 0), adjustment neither helps nor
hurts, as theory predicts: power is 71.6\% unadjusted versus 70.4\%
adjusted, and the type-I error is 2.5\% and 2.9\% respectively (2.2\%
and 3.2\% at full strength)---all within Monte Carlo error of the
nominal 2.5\%. As prognostic strength grows to the trial-calibrated
value (SD of the prognostic index $0.416$ on the log-hazard scale),
adjusted power rises to 80.9\% against 72.6\% unadjusted---an
efficiency gain equivalent to roughly 70 additional events (about a
23\% larger event count), obtained
at zero cost in bias or type-I error. This is the quantitative case
behind current regulatory encouragement of pre-specified covariate
adjustment~\cite{fda2023,kahan2014}.

\begin{table}[htbp]
\centering\small
\caption{Simulation Study A: covariate-adjustment efficiency (1{,}000
replicates per scenario, $n=322$, one-sided $\alpha=0.025$). Power at
HR $=0.72$; type-I error at HR $=1$. ``Bias'' contrasts each estimator
with the conditional log-HR used to generate the data.}
\label{tab:simA}
\begin{tabular*}{\textwidth}{@{\extracolsep{\fill}} l l c c c c}
\toprule
Prognostic strength & HR & Power (unadj.) & Power (adj.) & Bias (unadj.) & Bias (adj.) \\
\midrule
None (SD $0.000$)     & 1.00 & 0.025 & 0.029 & $-0.006$ & $-0.006$ \\
Moderate (SD $0.208$) & 1.00 & 0.025 & 0.025 & $+0.004$ & $+0.005$ \\
Full (SD $0.416$)     & 1.00 & 0.022 & 0.032 & $+0.000$ & $-0.001$ \\
None (SD $0.000$)     & 0.72 & 0.716 & 0.704 & $+0.000$ & $-0.009$ \\
Moderate (SD $0.208$) & 0.72 & 0.725 & 0.766 & $+0.012$ & $-0.006$ \\
Full (SD $0.416$)     & 0.72 & 0.726 & \textbf{0.809} & $+0.033$ & $-0.010$ \\
\bottomrule
\end{tabular*}
\end{table}

\begin{figure}[htbp]
  \centering
  \includegraphics[width=0.85\textwidth, height=0.45\textheight, keepaspectratio]{fig_sim_efficiency.pdf}
  \caption{Simulation Study A: empirical power of the unadjusted versus
  covariate-adjusted Cox test as a function of prognostic strength
  (target HR $0.72$, $n = 322$, 1{,}000 replicates per scenario).}
  \label{fig:simA}
\end{figure}

The bias column also teaches a subtler lesson: at full prognostic
strength, the unadjusted estimator averages $+0.033$ on the log scale
relative to the data-generating conditional hazard ratio---that is, an
unadjusted marginal HR of about $0.745$ when the conditional truth is
$0.72$. This attenuation toward the null is the hazard ratio's
non-collapsibility~\cite{martinussen2013}: unlike a linear-model
coefficient or a risk difference, the marginal HR changes when strong
prognostic covariates are removed from the model even without any
confounding. Consequently, ``the unadjusted and adjusted HRs estimate
different quantities'' is not a pedantic caveat but a measurable
effect, and it explains why adjusted and unadjusted estimates in
randomized trials need not coincide (and why the near-coincidence of
$0.623$ and $0.621$ in our real data reflects the modest prognostic
strength of the adjustment set rather than a mathematical necessity).

\subsection{Study B: operating characteristics of the sequential design}

Simulating the full monitoring procedure (3{,}000 trials per scenario,
event-triggered looks with calendar-time censoring) validates the
design and quantifies the value of the futility boundary
(Table~\ref{tab:simB}, Figure~\ref{fig:simB}). Under the alternative
(HR $0.65$), the design rejects in 82.9\% of trials---close to the
design-theoretic 90\%, with the shortfall attributable to the
discreteness of event-triggered looks and the normal approximation
in the boundary crossing at the interim snapshots. Stopping is
dominated by the second interim (42.2\%) rather than the first (18.6\%),
and the expected event count is 183.4 versus 235 at completion: the
design saves, on average, 51.6 events (roughly 22\%) when the drug
works. Under the null with futility adhered to, the type-I error is
2.4\% and the trial stops for futility at the first interim in 54.9\%
of cases---the aggressive Hwang--Shih--DeCani spending doing exactly
what it was designed to do. Ignoring futility (the non-binding option)
inflates the type-I error only trivially, to 2.7\%, confirming the
design's stated robustness.

\begin{table}[htbp]
\centering\footnotesize
\caption{Simulation Study B: simulated operating characteristics of the
group-sequential design (3{,}000 trials per scenario).}
\label{tab:simB}
\begin{tabular*}{\textwidth}{@{\extracolsep{\fill}} l c c c c c c}
\toprule
Scenario & Eff.\ IA1 & Eff.\ IA2 & Final & Total eff. & Exp.\ events & Events saved \\
\midrule
HR $=0.65$ (as designed)   & 0.186 & 0.422 & 0.222 & \textbf{0.829} & 183.4 & 51.6 \\
HR $=1$, futility adhered  & 0.002 & 0.008 & 0.014 & 0.024 & 152.7 & 82.4 \\
HR $=1$, futility ignored  & 0.001 & 0.009 & 0.017 & 0.027 & 234.4 & 0.6 \\
\bottomrule
\end{tabular*}
\end{table}

\begin{figure}[htbp]
  \centering
  \includegraphics[width=0.85\textwidth, height=0.45\textheight, keepaspectratio]{fig_sim_gs_oc.pdf}
  \caption{Simulation Study B: simulated stopping probabilities by
  analysis and reason under the alternative (right) and under the null
  with futility adherence (left).}
  \label{fig:simB}
\end{figure}

\subsection{Study C: interval estimation under non-proportional hazards}

When the true effect is time-varying (HR $0.45$ in year one, weakening
to $0.80$ thereafter) and the analysis fits a single proportional-
hazards Cox model, the estimator targets a weighted average of the two
regimes (simulation estimand: log-HR $= -0.407$, i.e.\ a pooled HR of
about $0.666$). Over 500 replicates, the Wald interval achieved 92.6\%
coverage of this estimand and the percentile bootstrap 90.4\%---both
marginally below their nominal 95\% (Table~\ref{tab:simC}). The
bootstrap is slightly narrower (mean width $0.434$ vs.\ $0.443$) but
not more accurate. The lesson is double-edged: neither standard
interval collapses under moderate non-proportional hazards---the
practical damage is a shifted estimand rather than drastically wrong
uncertainty---yet the modest undercoverage confirms that
misspecification costs are real, reinforcing the value of the
diagnostics and sensitivity analyses performed in
Section~\ref{sec:res-tx}.

\begin{table}[htbp]
\centering\footnotesize
\caption{Simulation Study C: 95\% interval performance under a
time-varying hazard ratio (500 replicates, $n=335$; 100 bootstrap
resamples per replicate).}
\label{tab:simC}
\begin{tabular*}{0.80\textwidth}{@{\extracolsep{\fill}} l c c c}
\toprule
Method & Coverage & Mean width & Mean log-HR \\
\midrule
Cox Wald CI            & 0.926 & 0.443 & $-0.407$ \\
Bootstrap percentile & 0.904 & 0.434 & $-0.407$ \\
\bottomrule
\end{tabular*}
\end{table}

\begin{figure}[htbp]
  \centering
  \includegraphics[width=0.85\textwidth, height=0.45\textheight, keepaspectratio]{fig_sim_coverage.pdf}
  \caption{Simulation Study C: Wald versus bootstrap percentile
  intervals for the log hazard ratio across 500 simulated trials under
  a time-varying true effect. The dashed line marks the simulation
  estimand.}
  \label{fig:simC}
\end{figure}
